% ID: SEL-D0039-OB-E04
% ESTADO: cotejada. Segunda pasada visual y borrador compilado revisados.
% Fuente: Selectividad/matematicas_aplicadas/Ua4MATEMÁTICAS APLICADAS examen.pdf
% SHA-256: 78ff6bdd48cbac71e6cc47d6ddf3dcb9492540ecb9cbe2032ebfec02e39fb66d
% Página: 2; opción B; ejercicio original 4.

Un fabricante de tuberías de PVC sabe que la distribución de los diámetros interiores de los tubos de conducción de agua que produce sigue una ley Normal con varianza $\sigma^2=0.25\,mm^2$. Para estimar el diámetro medio de esas tuberías, toma una muestra aleatoria de 64 tubos y comprueba que el diámetro medio de esa muestra es de 20 mm.
\begin{enumerate}[label=\alph*)]
\item \textbf{(1.5 puntos)} Calcule un intervalo de confianza, con un nivel del 98\%, para la media de los diámetros de los tubos que fabrica.
\item \textbf{(1 punto)} Halle el tamaño mínimo que debe tener una muestra de esa distribución para que la amplitud de un intervalo de confianza, con ese mismo nivel de confianza, sea inferior a 2 mm.
\end{enumerate}
